%===========================================================
% Capítulo: Resultados (I) — Protocolo de evaluación, métricas y estimación de la incertidumbre
%===========================================================
\chapter{Resultados}
\label{chap:resultados}

\section*{Protocolo de evaluación a nivel de sujeto}
Todos los resultados se estiman a nivel de \emph{sujeto/slide}, porque la decisión clínica de cribado ocurre en esa granularidad y porque la correlación intra–slide entre parches puede inflar artificialmente el desempeño si se evalúa parche a parche. Para cada pliegue externo del esquema anidado, el modelo y sus hiperparámetros se seleccionan exclusivamente con datos del entrenamiento interno y, una vez fijados, generan puntuaciones de probabilidad en el conjunto de prueba externo, primero a nivel de parche y luego agregadas a nivel de sujeto mediante la funcional $\mathcal{A}\{\cdot\}$ seleccionada (media, logit–media o softmax-$\beta$). Las predicciones \emph{out-of-fold} concatenadas sobre los $K_\text{out}$ pliegues constituyen una muestra imparcial $\{(\hat{p}_s,y_s)\}_{s=1}^{S}$, con $\hat{p}_s\in[0,1]$ la probabilidad calibrada para el sujeto $s$ y $y_s\in\{0,1\}$ su etiqueta. Todo lo que sigue —curvas ROC/PR, calibración, umbrales y análisis de costos— se deriva de esa muestra, lo que garantiza que ninguna estimación reutiliza datos para selección y evaluación simultáneamente \cite{Tsamardinos2022Patterns}.

\section*{AUC-ROC y AUPRC con bandas de confianza}
La capacidad discriminativa global se resume por el área bajo la curva ROC (AUC-ROC) y, dado que el problema de cribado puede presentar prevalencia baja, por el área bajo la curva de precisión–recobrado (AUPRC), que es más sensible a los verdaderos positivos cuando la clase positiva es minoritaria. Sea $R(\tau)=(\mathrm{FPR}(\tau),\mathrm{TPR}(\tau))$ el punto ROC al variar el umbral $\tau$ sobre las probabilidades $\hat{p}_s$; la AUC-ROC se estima por integración trapezoidal sobre la envolvente empírica y su incertidumbre se cuantifica con \emph{bootstrap} estratificado por sujeto: se re–muestran $S$ pares $(\hat{p}_s,y_s)$ con reemplazo $B$ veces (típicamente $B=2000$), se reconstruyen las curvas y se reporta el IC95\,\% (percentil o BCa) de las áreas. En paralelo, para AUPRC se construye la curva $\big(\mathrm{Recall}(\tau),\mathrm{Precision}(\tau)\big)$ y se integra de igual manera; dado su carácter no lineal y su sensibilidad a la prevalencia, también se informa la precisión de referencia $\pi/(1-\pi)$ correspondiente a un clasificador aleatorio con la misma prevalencia $\pi=\frac{1}{S}\sum_s y_s$, lo que contextualiza el valor absoluto de AUPRC en cada cohorte. El uso de \emph{bootstrap} a nivel de sujeto, y nunca a nivel de parche, es clave para que los intervalos capturen la variabilidad realmente relevante (inter–sujeto) y no la variabilidad espuria intra–slide.

\section*{Sensibilidad, especificidad y exactitud balanceada al punto operativo}
Aunque AUC-ROC y AUPRC resumen la discriminación global, la adopción clínica requiere fijar un \emph{punto operativo} que refleje los costos relativos de falsos negativos y falsos positivos. Definimos dos puntos: el de \emph{Youden} $\tau_J=\arg\max_\tau\{\mathrm{TPR}(\tau)-\mathrm{FPR}(\tau)\}$ y el umbral de \emph{Bayes} $\tau^\ast=C_{\mathrm{FP}}/(C_{\mathrm{FP}}+C_{\mathrm{FN}})$ cuando existe una matriz de costos explícita (ver Cap.~\ref{chap:metodologia}). En ambos se reportan sensibilidad $\mathrm{TPR}$, especificidad $\mathrm{TNR}$ y exactitud balanceada $\mathrm{ACC}_b=\tfrac{1}{2}(\mathrm{TPR}+\mathrm{TNR})$ junto con IC95\,\% calculados mediante intervalos de Wilson para proporciones o por \emph{bootstrap} de sujetos, lo que evita el sesgo de Wald en tamaños moderados. Estas estimaciones se acompañan de PPV y NPV contextualizados por la prevalencia de la cohorte; como guía interpretativa, si $\hat{p}=\Pr(y{=}1|x)$ está bien calibrada, entonces $\mathrm{PPV}(\tau)=\frac{\pi\,\mathrm{TPR}(\tau)}{\pi\,\mathrm{TPR}(\tau)+(1-\pi)\,\mathrm{FPR}(\tau)}$ y $\mathrm{NPV}(\tau)=\frac{(1-\pi)\,\mathrm{TNR}(\tau)}{(1-\pi)\,\mathrm{TNR}(\tau)+\pi\,\mathrm{FNR}(\tau)}$, de modo que mejoras marginales en especificidad pueden traducirse en incrementos sustanciales de PPV cuando $\pi$ es baja, lo que justifica explorar el frente de Pareto sensibilidad–especificidad más allá de un único umbral.

\section*{Comparaciones entre modelos y pruebas pareadas}
Para contrastar clasificadores con igual particionado externo se usan pruebas \emph{pareadas} sobre las mismas unidades de evaluación. La comparación de AUC-ROC puede abordarse con el test no paramétrico de DeLong, que modela la varianza de las áreas como función de los vectores de indicadores de orden para pares positivos/negativos, o bien, de forma más robusta cuando las distribuciones son irregulares, con un \emph{bootstrap} pareado de sujetos re–calculando la diferencia de áreas $\Delta\mathrm{AUC}$ en cada réplica y reportando su IC95\,\%; si el intervalo no cruza cero se considera evidencia de superioridad. Para métricas a umbral fijo (sensibilidad, especificidad, exactitud), la prueba de McNemar sobre la matriz $2\times 2$ de desacuerdos a nivel de sujeto cuantifica si la diferencia de errores es significativa; cuando los recuentos son bajos, la corrección de continuidad o el \emph{exacto binomial} son preferibles. Con múltiples comparaciones entre familias de rasgos (LBP vs.\ GLCM vs.\ FFT y combinaciones), se controla la FDR para mantener un error de descubrimiento aceptable a nivel estudio. El tamaño del efecto se expresa como diferencia absoluta de AUC o como \emph{odds ratio} de aciertos relativos al umbral clínico, ambos con intervalos de confianza, lo que evita conclusiones basadas sólo en $p$–valores.

\section*{Curvas ROC/PR con bandas y análisis de dominancia}
Las figuras incluyen curvas ROC y PR promediadas sobre los pliegues externos pero trazadas a partir de las predicciones \emph{out-of-fold} concatenadas, lo que impide la doble contabilidad de muestras y permite dibujar bandas de confianza por \emph{bootstrap} a lo largo de toda la curva, no sólo en el área. Esta visualización por bandas es útil para identificar regiones de operación donde un modelo domina a otro de manera \emph{uniforme} (curva y banda entera por encima), un criterio más exigente que la sola superioridad de AUC. En la curva PR se señala la línea de referencia de precisión $\pi$ y se resaltan los puntos operativos $\tau_J$ y $\tau^\ast$ con sus barras de incertidumbre vertical y horizontal, facilitando una lectura clínica inmediata de los compromisos sensibilidad–PPV.

\section*{Calibración, confiabilidad y decisión clínica}
La calibración se evalúa con curvas de confiabilidad construidas en cuantiles de probabilidad (p.\,ej., 10 bins equiprobables), informando el \emph{Brier score} total y sus componentes de \emph{calibración} y \emph{refinamiento}. Un calibrador isotónico ajustado únicamente sobre predicciones \emph{out-of-fold} del entrenamiento externo y aplicado después a los sujetos del pliegue correspondiente mejora la correspondencia probabilística sin contaminar la evaluación; el intercepto de calibración (idealmente cercano a 0) y la pendiente (idealmente 1) se estiman por regresión logística de $y$ sobre $\logit(\hat{p})$. Dado un umbral de implementación, las curvas de \emph{net benefit} de la \emph{Decision Curve Analysis} (DCA) permiten comparar el beneficio clínico neto del modelo frente a \emph{tratar a todos} o \emph{tratar a ninguno} en el rango de probabilidades de decisión relevantes; cuando el modelo arroja net benefit superior de forma sostenida, ello respalda la adopción operativa más allá de la discriminación pura \cite{Vickers2006MedDec}.

\section*{Coherencia con la física: verificaciones observacionales}
Más allá de la métrica global, se verifica que los rasgos que impulsan la decisión están alineados con el modelo físico de difracción. La media y los intervalos de los picos radiales $(k_1,k_2)$ en la clase positiva se desplazan hacia valores más altos respecto a la clase negativa; la razón espectral $R_{HL}$ crece de forma consistente con franjas más juntas y estructura de alta frecuencia más rica en AF; la anisotropía angular aumenta y, cuando se proyecta en modos periódicos, el componente $m{=}2$ gana peso indicando elipticidad del patrón. En textura, la entropía de LBP y el contraste de GLCM son mayores en positivos, mientras que la homogeneidad decrece, una triada compatible con la ruptura de circularidad y el enriquecimiento de micro–alternancias esperadas en morfologías elongadas. Estas observaciones se reportan con diferencias de medias y IC95\,\% por \emph{bootstrap} de sujetos, y no sólo con pruebas de significancia, para reforzar la interpretación causal–mecánica y no meramente correlacional \cite{Kiani2022TextureSurvey,Chen2022Photonics}.

\section*{Análisis de errores, estratificación por SNR y robustez}
La tasa de falsos negativos se descompone por cuartiles de $\mathrm{SNR}_B$ definida en el dominio espectral: si los FN se concentran en el cuartil inferior, se documenta explícitamente el deterioro de sensibilidad y se exploran mitigaciones (umbral más bajo en condiciones de SNR baja o descarte automático por QC con repetición de adquisición). Se realiza la misma estratificación por métricas de iluminación global (media y varianza de intensidad) y por proxies instrumentales cuando están disponibles (tamaño de píxel $p$ o distancia estimada $z$); la estabilidad de AUC y de la curva PR a través de estratos apoya la transferencia del sistema entre configuraciones LDHM razonablemente similares. Como verificación adicional, se perturban de manera controlada las imágenes de prueba con ruido aditivo y cambios suaves de ganancia para cuantificar el gradiente de degradación de AUC; una pendiente pequeña en estas perturbaciones indica que los descriptores están capturando estructura interferométrica robusta y no artefactos fotométricos superficiales.

\section*{Interpretabilidad cuantitativa y validación cruzada de explicaciones}
Los diagramas \emph{beeswarm} de SHAP y los \emph{waterfall} por sujeto confirman que los rasgos dominantes en la predicción positiva son $k_1$, $R_{HL}$, la pendiente espectral menos negativa, la entropía/contraste LBP y el contraste/entropía de GLCM; la homogeneidad elevada y la anisotropía baja empujan hacia negativo. Para evitar conclusiones sobreajustadas a un pliegue, las explicaciones se reportan por pliegue externo y se promedian únicamente a nivel de ranking de rasgos, no de valores absolutos, ya que los valores de Shapley dependen de la distribución de rasgos en cada entrenamiento. Como control de sanidad, se repite el análisis con permutación de etiquetas; la desaparición de patrones estructurados en las explicaciones bajo esta permutación respalda que las señales explotadas no emergen de sesgos de particionado.

\section*{Reproducibilidad, potencia y límites de inferencia}
Se documentan las semillas aleatorias y versiones de librerías de cada corrida, junto con las predicciones \emph{out-of-fold} y los calibradores por pliegue, de modo que cualquier tercero pueda reproducir cada figura y cada tabla a partir de artefactos almacenados. En cuanto a potencia estadística, aun cuando la AUC supere el umbral objetivo, los intervalos de confianza anchos motivan cautela: un ancho de IC por encima de 0.15 sugiere que el número de sujetos es todavía insuficiente para afirmaciones concluyentes y que los resultados deben interpretarse como evidencia de viabilidad más que de validación clínica. En esta línea, se explicita que el estudio carece por ahora de una cohorte externa independiente; los buenos resultados internos, al estar sustentados por validación anidada, son alentadores pero requieren confirmación \emph{out-of-domain} para reclamar generalización fuerte.

\section*{Síntesis de la primera parte de resultados}
El sistema demuestra discriminación alta medida por AUC-ROC y AUPRC, con sensibilidad elevada en puntos operativos que priorizan el costo del falso negativo, y produce probabilidades bien calibradas que se traducen en beneficios clínicos netos en rangos de decisión relevantes. Las señales dominantes coinciden con lo que predice la física de la difracción en LDHM para células elongadas, y la robustez frente a SNR variable y pequeñas perturbaciones fotométricas sugiere que el modelo captura estructura interferométrica y no sólo texturas superficiales. Estos hallazgos, anclados en estimación honesta con validación anidada y cuantificación explícita de la incertidumbre, constituyen una base sólida para el análisis comparativo entre familias de rasgos y para la discusión crítica que seguirá en la segunda parte de este capítulo.

%-----------------------------------------------------------
% En el siguiente bloque (Resultados II) se presentarán comparativas entre familias de rasgos,
% ablasiones, análisis por subgrupos, decisión clínica con DCA y un estudio detallado de casos.
%-----------------------------------------------------------


%===========================================================
% Capítulo: Resultados (II) — Comparativas entre familias de rasgos, ablasiones y análisis por subgrupos
%===========================================================

\section*{Comparativas entre familias de rasgos sin reconstrucción}
La comparación sistemática entre familias de rasgos confirma que la señal discriminativa más estable emerge cuando se preserva la triada físico–estadística completa (LBP + GLCM + FFT/anisotropía) y se evalúa estrictamente a nivel de sujeto. Cuando se analiza cada familia por separado, el perfil espectral radial junto con los indicadores de anisotropía angular tiende a capturar de manera directa los cambios de escala y direccionalidad de las franjas asociados a elongación morfológica, traduciéndose en curvas ROC que se elevan de forma uniforme en regiones de baja tasa de falsos positivos; por su parte, la textura local (LBP) y la co–ocurrencia (GLCM) reflejan el enriquecimiento del microcontraste interferométrico y su irregularidad de segundo orden, lo que favorece la precisión en rangos de alta sensibilidad al desplazar la curva PR hacia mayores valores de precisión a igual recobrado. La combinación de las tres familias produce una envolvente que, sin recurrir a reconstrucción numérica, sintetiza escala, periodicidad y anisotropía en un espacio de baja dimensión con clara correspondencia física; esta sinergia se manifiesta en mejoras modestas pero consistentes en AUC-ROC y, sobre todo, en AUPRC cuando la prevalencia efectiva es baja, porque la energía de alta frecuencia y la anisotropía aportan separabilidad temprana en el ranking, mientras que LBP/GLCM consolidan la robustez en presencia de ruido y variaciones de iluminación \cite{Kiani2022TextureSurvey,Chen2022Photonics}.

\section*{Ablaciones y parquedad con interpretabilidad física}
Las ablasiones confirman que el beneficio incremental de añadir rasgos debe ponderarse contra el aumento de varianza estimativa en regímenes de pocos sujetos. En esta dirección, la selección con FDR seguida de importancias por permutación y RFE identifica invariablemente un subespacio compacto dominado por el primer pico radial $k_1$, la razón espectral $R_{HL}$, la pendiente en banda media y dos o tres descriptores de textura (entropía y contraste de LBP y GLCM), a los que se suma, como componente angular, un modo $m{=}2$ que captura elipticidad del patrón. Un hallazgo consistente es que recortar el vector de rasgos a este núcleo altera muy poco la discriminación global y mejora la estabilidad inter–pliegue, a la vez que facilita explicaciones \emph{post-hoc} con significado físico directo: mayores $k_1$ y $R_{HL}$ inclinan la predicción hacia AF porque codifican franjas más juntas y contenido direccional más marcado; valores elevados de entropía y contraste reflejan el desorden periódico inducido por la pérdida de circularidad y la rigidez celular. Este equilibrio entre parquedad e interpretabilidad es crucial para auditar el comportamiento del sistema, pues permite mapear, pliegue a pliegue, qué magnitudes gobiernan la decisión y verificar su coherencia con el modelo de difracción de Fresnel sin necesidad de reconstruir amplitud o fase.

\section*{Análisis por subgrupos: SNR espectral, fotometría e instrumentación}
La robustez se evalúa estratificando a los sujetos por cuartiles de $\mathrm{SNR}_B$, definida en torno al primer lóbulo espectral, y por métricas fotométricas globales (media y desviación estándar de intensidad después de percentiles y estandarización). La degradación de sensibilidad y la contracción de AUPRC que se observan en el cuartil de SNR más bajo mantienen, no obstante, la forma de la curva ROC con pendientes iniciales pronunciadas, atributo característico de modelos que extraen estructura interferométrica real; cuando la SNR mejora, la precisión se separa de la línea de referencia de forma más precoz, lo que se traduce en ganancias prácticas de PPV para umbrales fijos. Un análisis paralelo considera proxies instrumentales, como el tamaño de píxel efectivo $p$ o variaciones de distancia $z$ cuando están disponibles; al expresar las frecuencias en fracción de Nyquist y fijar bandas log–espaciadas, las curvas por estratos de $p$ y $z$ muestran perfiles muy próximos, lo que respalda que la normalización espectral amortigua desplazamientos de dominio suaves. Desde la perspectiva de despliegue, estos resultados recomiendan dos medidas de ingeniería de procesos: un chequeo automático de SNR por sujeto para informar al operador cuando una adquisición cae por debajo de umbral de calidad, y una política de repetición rápida de la captura en esos casos para recuperar sensibilidad sin necesidad de recalibrar el modelo.

\section*{Comparaciones pareadas y estabilidad inter–pliegue}
Las diferencias entre configuraciones —por ejemplo, FFT/anisotropía sola frente a triada completa— se evalúan con pruebas pareadas a nivel de sujeto. En la práctica, el \emph{bootstrap} pareado ofrece una lectura más transparente que el test de DeLong cuando las distribuciones de puntuaciones son asimétricas: al re–muestrear sujetos y recomputar la diferencia de áreas $\Delta\mathrm{AUC}$ o $\Delta\mathrm{AUPRC}$, los intervalos de confianza muestran si la ventaja de la combinación es robusta o si, por el contrario, se debe a una fluctuación particular del particionado. La estabilidad inter–pliegue se ilustra superponiendo curvas ROC/PR por pliegue y resaltando la banda de confianza de la envolvente global; cuando la triada de rasgos es dominante, las curvas individuales se apilan de manera más estrecha que en familias aisladas, lo que sugiere que los descriptores redundantes en física (escala, textura, direccionalidad) ejercen un efecto de \emph{ensamble estadístico} que amortigua la varianza.

\section*{Calibración y utilidad clínica bajo diferentes prevalencias}
Dado que la utilidad clínica depende de la prevalencia poblacional, se exploran escenarios hipotéticos de despliegue ajustando la base de probabilidad mediante la identidad de Bayes. Si $\pi$ es la prevalencia en la cohorte interna y $\pi'$ la prevalencia en el entorno de destino, la razón de verosimilitudes estimada por el modelo, $\mathrm{LR^+}=\mathrm{TPR}/\mathrm{FPR}$ y $\mathrm{LR^-}=\mathrm{FNR}/\mathrm{TNR}$ en el umbral elegido, permite actualizar las probabilidades previas: $\mathrm{post\ odds}=\mathrm{pre\ odds}\times \mathrm{LR}$, con $\mathrm{pre\ odds}=\pi'/(1-\pi')$. Este cálculo muestra cómo, incluso con TPR alta, la PPV puede variar notablemente al pasar de una clínica de alta sospecha a un tamizaje universal; por ello, las curvas de \emph{net benefit} de la \emph{Decision Curve Analysis} se reportan sobre un rango de probabilidades de decisión que refleja distintos perfiles de prevalencia, poniendo en evidencia que la misma discriminación puede tener beneficios netos diferentes según el escenario operativo \cite{Vickers2006MedDec}. La calibración isotónica, aplicada por pliegue externo sobre predicciones \emph{out-of-fold}, mantiene una pendiente cercana a la unidad y un intercepto pequeño, lo que justifica usar directamente las probabilidades para decisiones informadas por umbral de costo.

\section*{Coherencia de las explicaciones y validación cruzada de XAI}
La inspección de explicaciones con SHAP confirma que el núcleo parco de rasgos domina la decisión de manera consistente entre pliegues: $k_1$ y $R_{HL}$ elevan las \emph{log-odds} hacia la clase positiva cuando superan cuantiles medios, mientras que una homogeneidad GLCM alta ejerce el efecto contrario. Un aspecto importante es que las contribuciones absolutas no deben promediarse sin cuidado entre pliegues, pues dependen de la distribución de rasgos en el entrenamiento; para evitar este sesgo, se reporta el \emph{ranking} medio de importancia junto con bandas de cuantil por pliegue. Al permutar etiquetas dentro de cada entrenamiento externo, las explicaciones se desestructuran y el ranking se aplana, reforzando la hipótesis de que el modelo explota relaciones físicas reales y no artefactos de particionado. Adicionalmente, al proyectar las contribuciones sobre el eje de frecuencia radial, se aprecia un gradiente de influencia que crece desde la banda baja hacia la media–alta, en sintonía con la expectativa de que franjas más juntas (altas frecuencias relativas) aportan evidencia de morfología falciforme.

\section*{Errores característicos y estudio de casos}
El análisis cualitativo de falsos negativos frecuentes revela dos patrones característicos: hologramas con contraste micro–local insuficiente a pesar de la ecualización adaptativa, y patrones donde la estructura angular está presente pero distribuida en múltiples orientaciones débiles, lo que diluye la anisotropía agregada; en ambos casos, reducir el umbral operativo mejora la captura de positivos al coste de algunos falsos positivos adicionales, un intercambio aceptable en contextos de cribado. Entre los falsos positivos, destacan parches con micro–artefactos de borde o suciedad que imitan alternancias de alta frecuencia, lo que sugiere endurecer el filtro de SNR espectral o introducir un descriptor de \emph{coherencia radial} que penalice energía de alta frecuencia dispersa sin organización anular. Estos hallazgos alimentan un ciclo de mejora incremental del pipeline que respeta el principio de parquedad y se mantiene dentro del marco físico–estadístico propuesto.

\section*{Validez externa y límites de inferencia}
Aun cuando la discriminación interna sea elevada y la calibración adecuada, la validez externa requiere evaluar el modelo en una cohorte verdaderamente \emph{out-of-domain} con diferencias instrumentales y de preparación de muestra controladas. La normalización por Nyquist, el uso de bandas log–espaciadas y la multiescala en LBP/GLCM mitigan desplazamientos suaves, pero no sustituyen la verificación externa. En ausencia de esa cohorte, la incertidumbre se comunica explícitamente mediante IC amplios cuando el número de sujetos es moderado, y se propone una estrategia de \emph{congelar} el subespacio parco de rasgos y recalibrar únicamente el umbral o el calibrador isotónico con un pequeño lote de validación del nuevo sitio, lo que reduce el riesgo de sobreajuste y mantiene el anclaje físico de las magnitudes utilizadas.

%===========================================================
% Capítulo: Resultados (III) — Decisión clínica, utilidad y hojas de ruta de mejora
%===========================================================

\section*{Decisión clínica y curvas de utilidad}
La traducción de probabilidades calibradas a decisiones exige mapear explícitamente los costos relativos y la prevalencia. Cuando $C_{\mathrm{FN}}\gg C_{\mathrm{FP}}$, el umbral bayesiano $\tau^\ast=\frac{C_{\mathrm{FP}}}{C_{\mathrm{FP}}+C_{\mathrm{FN}}}$ cae en la región de alta sensibilidad, y la \emph{Decision Curve Analysis} muestra un beneficio neto superior al de estrategias triviales para un rango amplio de probabilidades de decisión. Esta superioridad se mantiene cuando la prevalencia desciende, siempre que la especificidad no se degrade por debajo de un nivel mínimo; por ello, la operacionalización recomienda un umbral primario de cribado y un umbral secundario más estricto para confirmar internamente antes de derivar a pruebas de referencia, articulando un flujo de trabajo de dos etapas que optimiza sensibilidad sin cargar excesivamente al laboratorio confirmatorio \cite{Vickers2006MedDec}. La documentación de estos umbrales, junto con matrices de confusión y PPV/NPV por escenario de prevalencia, proporciona a los clínicos una guía clara de interpretación y uso.

\section*{Hojas de ruta de mejora y compatibilidad con POCT}
Los resultados orientan dos líneas de mejora conservando la filosofía de no reconstrucción. Primero, incorporar un descriptor de \emph{coherencia anular} basado en la varianza radial normalizada dentro de cada banda, que penaliza energía alta dispersa no organizada en anillos; esta medida, físicamente motivada, discrimina artefactos granulares de patrones celulares reales. Segundo, explorar un clasificador \emph{híbrido} que, manteniendo el subespacio parco de rasgos físicos, añada una proyección aprendida de muy baja dimensión con regularización fuerte (p.\,ej., un \emph{autoencoder lineal} o una capa densa con norma nuclear), de modo que cualquier ganancia sea atribuible a dependencias de orden superior y no a memorias espurias; la validación anidada y los IC asegurarán que las mejoras no sean producto del azar. Estas mejoras son compatibles con la ejecución en CPU y con plataformas LDHM de bajo costo ya descritas, y no requieren modificar la cadena de adquisición ni estimar fase, por lo que preservan el presupuesto computacional objetivo \cite{Greenbaum2014STM,Chen2022Photonics}.

\section*{Síntesis}
El cuerpo de resultados respalda que el análisis directo de hologramas crudos, formalizado como una proyección físico–estadística parca y evaluado con protocolos honestos, ofrece discriminación alta, probabilidades confiables y beneficios clínicos netos en escenarios de tamizaje, con explicaciones que se alinean con la física de la difracción. La estabilidad por subgrupos de SNR y la relativa invariancia a la instrumentación, junto con la claridad de las ablasiones, configuran un caso sólido para escalar la validación y explorar una transición controlada hacia entornos POCT, en los que la portabilidad, el bajo costo y la rapidez de inferencia constituyen ventajas diferenciales del enfoque propuesto \cite{Tsamardinos2022Patterns,Kiani2022TextureSurvey,Vickers2006MedDec,Greenbaum2014STM,Chen2022Photonics}.
